%% kapa: exam=1 round=25 year=2014 period=1 form=A subject=civil_law subject_ko=민법 seq=16 number=16
%% source: Q-Net 감정평가사 기출문제 (cst00309) artlSeq=5101207 file=extracted/1008/5101207/2014년 제25회 감정평가사 1차시험 기출문제/1교시(A형).hwp
%% conversion: kordoc 4.24.51 -> tools/gen1.py (auto); review=proofread-opus-orchestrator; flags=none
\begin{question}{16}
불공정한 법률행위에 관한 설명으로 옳지 않은 것은? (다툼이 있으면 판례에 의함)
\begin{choices}
  \item 어업권 소멸로 인한 손실보상금의 분배에 관하여 어촌계 총회의 결의가 현저하게 불공정한 경우, 폭리행위가 될 수 있다.
  \item 매매계약이 불공정한 법률행위에 해당하는 경우, 특별한 사정이 없는 한, 불이익을 입게 되는 당사자가 그 불공정성을 이유로 제소를 하지 못하게 하는 부제소합의는 무효이다\uline{.}
  \item 매매계약에서 의사표시를 한 자가 궁박 상태에 있었더라도 상대방이 그와 같은 사정을 알면서 이를 이용하려는 의사가 없으면 그 계약은 불공정한 법률행위가 되지 않는다.
  \item 불공정한 법률행위를 이유로 무효를 주장하는 자는 피해자의 궁박, 경솔과 무경험을 모두 증명해야 한다.
  \item 불공정한 법률행위로서 무효인 법률행위는 추인에 의해서 유효로 될 수 없다.
\end{choices}
\end{question}
\answer{4}
